% Generated draft: requires mathematical and source review.
\documentclass[11pt]{article}
\usepackage[margin=25mm]{geometry}
\usepackage{fontspec}
\setmainfont{Latin Modern Roman}
\usepackage{amsmath,amssymb,mathtools,unicode-math}
\setmathfont{Latin Modern Math}
\usepackage{enumitem,needspace,xurl,hyperref}
\setlength{\parindent}{0pt}
\setlength{\parskip}{5pt}
\setlength{\emergencystretch}{3em}
\newcommand{\N}{\mathbb{N}}
\newcommand{\Z}{\mathbb{Z}}
\newcommand{\Q}{\mathbb{Q}}
\newcommand{\R}{\mathbb{R}}
\newcommand{\C}{\mathbb{C}}
\newcommand{\F}{\mathbb{F}}
\newcommand{\A}{\mathbb{A}}
\newcommand{\PP}{\mathbb{P}}
\newcommand{\E}{\mathbb{E}}
\newcommand{\eps}{\varepsilon}
\newcommand{\dd}{\,\mathrm{d}}
\begin{document}
\section*{UnsolvedMath ID: 30005068}
\textbf{Source problem code:} catalogue entry
\subsection{Definitions and mathematical statement}
The supplied record asks about boundary behavior of nonlocal minimal surfaces, with particular emphasis in its reconstructed statement on surfaces beyond planar graphs. It does not specify a kernel, dimension, admissible geometry, or notion of boundary equation. The following standard fractional-perimeter framework provides terminology only; its identification with the intended setting requires review.

Let \(n\geq 1\), \(s\in(0,1)\), and let \(\Omega\subset\R^n\) be open. For disjoint measurable sets \(A,B\subset\R^n\), write
\[
L_s(A,B)=\int_A\int_B\frac{1}{|x-y|^{n+s}}\,\dd y\,\dd x.
\]
For a measurable set \(E\subset\R^n\), put \(E^c=\R^n\setminus E\) and define its relative fractional perimeter by
\[
P_s(E;\Omega)=L_s(E\cap\Omega,E^c)+L_s(E\setminus\Omega,E^c\cap\Omega).
\]
In the finite-perimeter version of this framework, \(E\) is a minimizer with prescribed exterior data \(E_0\) if \(P_s(E;\Omega)<\infty\), \(E\setminus\Omega=E_0\setminus\Omega\) up to null sets, and
\[
P_s(E;\Omega)\leq P_s(F;\Omega)
\]
for every measurable \(F\subset\R^n\) having the same exterior data. These are explicit conventions for this explanatory framework, not hypotheses recovered from the source.

At a point \(x\) of a suitably chosen representative of \(\partial E\), define the fractional mean curvature, when the principal value exists, by
\[
H_s[E](x)=\operatorname{PV}\int_{\R^n}
\frac{\mathbf{1}_{E^c}(y)-\mathbf{1}_E(y)}{|x-y|^{n+s}}\,\dd y
=\lim_{\varepsilon\downarrow0}\int_{\R^n\setminus B_\varepsilon(x)}
\frac{\mathbf{1}_{E^c}(y)-\mathbf{1}_E(y)}{|x-y|^{n+s}}\,\dd y.
\]
The corresponding classical nonlocal mean-curvature equation is \(H_s[E](x)=0\).

\textbf{Requested problem, with unresolved qualifications.} Determine under which hypotheses the nonlocal mean-curvature equation for a nonlocal minimal surface continues to hold up to the boundary of its minimizing domain, including the setting described in the supplied reconstruction as beyond planar graphs. In the framework above, potential boundary contact points lie in \(\partial E\cap\partial\Omega\); the source does not specify which such points are intended, whether the equation means a pointwise principal-value identity there, an interior limiting identity, or a weak or viscosity formulation, or what boundary regularity is assumed.

Also characterize how generally boundary stickiness occurs. Here ``stickiness'' is only a descriptive term for attachment of a minimizing surface to the domain boundary or failure to attain prescribed exterior data continuously; the source supplies no precise predicate. It also gives no topology, measure, probability law, or parameter family defining ``typical.'' Thus a precise quantified assertion about typicality, and a universal boundary identity, cannot be extracted without additional hypotheses. The source presents research questions rather than a specified theorem to prove.
\subsection{Short English statement}
Determine when nonlocal minimal surfaces satisfy their mean-curvature equation up to the domain boundary, especially beyond planar graphs, and characterize the prevalence of boundary stickiness. The supplied source leaves the admissible surfaces, boundary formulation, definition of stickiness, and meaning of typicality unspecified. The fractional-perimeter definitions above provide a conditional framework for review.
\subsection{Sources}
\begin{itemize}
\item Local registry: \texttt{data/unsolvedmath/problems.json}, with \texttt{catalog/problems.json} as the supplementary project registry; supplied permanent UnsolvedMath ID \(30005068\), title ``Boundary Curvature and Stickiness of Nonlocal Minimal Surfaces,'' source problem code \texttt{catalogue entry}. These identifiers refer to the supplied record; the registry files have not been independently inspected.
\item Supplied report citation: \emph{MFO-RIMS Tandem Workshop: Nonlocality in Analysis, Probability and Statistics (2023)}, Oberwolfach Reports, DOI \texttt{10.4171/owr/2022/15}, pp.\ 709--715. \url{https://doi.org/10.4171/owr/2022/15}. The literature-source label is ``Oberwolfach Report 15/2022.'' The supplied proposed year and year tag are 2022, whereas the citation names 2023; this discrepancy is unresolved.
\item Supplied original attribution: ``Nonlocal minimal surfaces,'' Serena Dipierro, joint work with Ovidiu Savin and Enrico Valdinoci. The original excerpt asks whether the nonlocal mean-curvature equation is satisfied up to the boundary and how typical stickiness is. No separate proposer is identified.
\item Supplied literature source: Dipierro--Savin--Valdinoci, ``boundary continuity'' (descriptive label supplied by the registry). \url{https://arxiv.org/abs/2305.15271}. No fuller bibliographic identification is supplied.
\item Literature assessment dated exactly \texttt{2026-08-22}: the supplied curation labels the problem ``partially solved'' and states that boundary continuity and stickiness have been rigorously classified in several smooth and singular geometries, while the general boundary-equation and typicality question in the extracted form remains open. This is supplied dated curation, not an independently verified literature assessment. The record further labels the statement \texttt{corrected\_verified} and reports reconstruction and checking against the official report; those verification claims are supplied metadata, not verification performed for this draft.
\end{itemize}
\end{document}
